\usepackage{piton}
\usepackage{markdown}
\usepackage{attachfile2}

\newtcbtheorem[number within=chapter]{CodeExampleBox}{\Lang{Codebeispiel}{Code example}}
  {scriptbox,breakable,borderline west={1.2pt}{0pt}{THMBlack}}{code}

\newfontfamily\ScriptPythonFont{DejaVuSansMono.ttf}[
  BoldFont=DejaVuSansMono-Bold.ttf,
  ItalicFont=DejaVuSansMono-Oblique.ttf,
  BoldItalicFont=DejaVuSansMono-BoldOblique.ttf
]
\SetPitonStyle[Python]{
  Keyword=\color[HTML]{234780}\bfseries,
  Keyword.Governing=\color[HTML]{234780}\bfseries,
  Keyword.Constant=\color[HTML]{234780}\bfseries,
  Operator.Word=\color[HTML]{234780}\bfseries,
  Name.Function=\color[HTML]{5A347A},
  Name.Namespace=\color[HTML]{285E61},
  Name.Builtin=\color[HTML]{285E61},
  Comment=\color[HTML]{42654F}\itshape,
  String=\color[HTML]{8B3B25},
  String.Doc=\color[HTML]{8B3B25}\itshape,
  Number=\color[HTML]{6C4679}
}

\ExplSyntaxOn
\NewDocumentCommand{\ScriptPythonListing}{m}{
  \group_begin:
    % Write indentation as space glyphs so it is part of the PDF text.
    \cs_set_protected:Npn \__piton_leading_space: {
      \int_gincr:N \g__piton_indentation_int
      \char_generate:nn { 32 } { 12 }
    }
    \cs_set_protected:Npn \__piton_tab: {
      \int_gadd:Nn \g__piton_indentation_int \l__piton_tab_size_int
      \prg_replicate:nn { \l__piton_tab_size_int } { \char_generate:nn { 32 } { 12 } }
    }
    \PitonInputFile[language=Python,font-command={\ScriptPythonFont\footnotesize},
      left-margin=0.8em,tab-size=8,break-lines,splittable]{#1}
  \group_end:
}
\ExplSyntaxOff

\newcommand{\ScriptCellLabelFont}{\normalfont\sffamily\fontsize{6pt}{7pt}\selectfont}
\NewDocumentCommand{\ScriptPythonSource}{m m}{%
  {\ScriptCellLabelFont
    \textattachfile[ucfilespec={#2},mimetype=text/x-python,color={black!55},
      description={\Lang{Python-Quellcode}{Python source}}]{#1}{Python}}%
}
\NewDocumentCommand{\ScriptPythonFile}{m}{%
  \ScriptPythonListing{#1}%
  \par\noindent\hfill\ScriptPythonSource{#1}{\ScriptPythonListingTitle}\par
}

\newtcolorbox{ScriptNotebookCell}[1]{%
  enhanced,frame hidden,colback=LectureBoxBackground,boxrule=0.25pt,
  borderline={0.25pt}{0pt}{black!55},
  sharp corners,boxsep=0pt,left=0pt,right=0.8em,top=3pt,bottom=10pt,
  before skip=5pt,after skip=5pt,
  overlay={
    \node[anchor=south east,inner sep=0pt,text=black!55,
      font=\ScriptCellLabelFont]
      at ([xshift=-3pt,yshift=2.5pt]frame.south east) {#1};
  }
}
\NewDocumentCommand{\ScriptNotebookCodeCell}{m m}{%
  \begin{ScriptNotebookCell}{\ScriptPythonSource{#1}{#2}}
    \ScriptPythonListing{#1}%
  \end{ScriptNotebookCell}%
}

% Notebook headings are local to the box and do not change script numbering.
\NewDocumentCommand{\ScriptNotebookHeading}{m}{\par\noindent\textbf{#1}\par}
\NewDocumentCommand{\ScriptNotebookImage}{m m m m}{\emph{#1}}
\NewDocumentCommand{\ScriptNotebookMarkdownCell}{m}{%
  \begin{ScriptNotebookCell}{Markdown}
    \small
    \leftskip=0.8em
    \let\markdownRendererHeadingOne\ScriptNotebookHeading
    \let\markdownRendererHeadingTwo\ScriptNotebookHeading
    \let\markdownRendererHeadingThree\ScriptNotebookHeading
    \let\markdownRendererHeadingFour\ScriptNotebookHeading
    \let\markdownRendererHeadingFive\ScriptNotebookHeading
    \let\markdownRendererHeadingSix\ScriptNotebookHeading
    \let\markdownRendererImage\ScriptNotebookImage
    \markdownInput[hybrid=false,html=false,fencedCode=true]{#1}%
    \par
  \end{ScriptNotebookCell}%
}

% Source paths are relative to the document root in all build modes.
\NewDocumentCommand{\PythonCode}{o m}{%
  \begingroup
    \directlua{
      if not script_python_listing then
        kpse.record_input_file("../_shared/commands/python-code.lua")
        script_python_listing = dofile("../_shared/commands/python-code.lua")
      end
      local path = "\luaescapestring{#2}"
      token.set_macro("ScriptPythonListingTitle", path:match("[^/]+$") or path)
    }%
    \begin{CodeExampleBox}{\IfNoValueTF{#1}{%
      \texttt{\expandafter\detokenize\expandafter{\ScriptPythonListingTitle}}}{#1}}{}
      \directlua{script_python_listing("\luaescapestring{#2}")}%
    \end{CodeExampleBox}%
  \endgroup
}
